\documentclass[12pt]{article}
\usepackage[utf8]{inputenc}
\usepackage[T1]{fontenc}
\usepackage{hyperref}
\title{Yemba Language Resources for K-12 Curriculum Development}
\author{Compiled by Claude Sonnet 5.5 (claude-sonnet-5-5) for Andr\'as Kornai}
\date{}
\begin{document}
\maketitle

\section{Introduction}
Yemba (ISO 639-3 ybb; Glottocode yemb1246) is a Grassfields Bantu language of the Menoua area around Dschang, West Region, Cameroon. The reference description says it is used as a language of instruction in some schools. The search results did not confirm this.

\section{Primary Educational Resources}
No textbooks or school programmes were found. The yemba.net site offers quizzes and tests alongside its dictionary.
\begin{itemize}
\item yemba.net, which has learning tools. \url{https://yemba.net/}
\end{itemize}

\section{Linguistic and Reference Materials}
\begin{itemize}
\item Online Yemba--French dictionary with over 3000 entries. It builds on the \emph{Petit Dictionnaire Y\'emba--Fran\c{c}ais} (Bird and Tadadjeu, CELY 1997, published by ANACLAC). \url{https://yemba.net/wdict/index.php?lang=en}
\item Webonary Yemba dictionary and overview. \url{https://www.webonary.work/yamba/overview/introduction/}
\item LDC, Grassfields Bantu Fieldwork: Dschang Lexicon, with sound files. \url{https://catalog.ldc.upenn.edu/LDC2003L01}
\item \emph{Analyse contrastive des parlers Y\'emba du d\'epartement de la M\'enoua}. \url{https://www.koeppe.de/titel_print_analyse-contrastive-des-parlers-yemba-du-departement-de-la-menoua-de-l-ouest-cameroun}
\item Yemba language (Wikipedia). \url{https://en.wikipedia.org/wiki/Yemba_language}
\item Omniglot page on Yemba writing. \url{https://omniglot.com/writing/yemba.htm}
\end{itemize}

\section{Religious and Community Resources}
The SIL Cameroon archive lists Yemba items. Its contents were not examined. \url{https://www.silcam.org/fr/resources/search/publisher/17604/worktype/15888/language/ybb}

\section{Digital and Audio Resources}
The LDC lexicon has audio for each entry. The yemba.net dictionary is also online. A font page for the Yemba Latin script exists. \url{https://www.fontspace.com/languages/info/ybb-latn}

\section{Curriculum Development Context}
Yemba has an online dictionary, a lexicon with audio, and an orthography associated with the CELY and ANACLAC committees. No evidence of classroom use was found. A curriculum would need primers and readers built on the existing dictionary.

\end{document}